\documentclass[11pt,a4paper]{article}
\usepackage[margin=1in]{geometry}
\usepackage{hyperref}
\usepackage{listings}
\usepackage{booktabs}
\usepackage{graphicx}
\usepackage{xcolor}
\usepackage{fancyhdr}
\usepackage{tocloft}

% Code listing style
\definecolor{codegreen}{rgb}{0,0.6,0}
\definecolor{codegray}{rgb}{0.5,0.5,0.5}
\definecolor{codepurple}{rgb}{0.58,0,0.82}
\definecolor{backcolour}{rgb}{0.95,0.95,0.92}

\lstdefinestyle{mystyle}{
    backgroundcolor=\color{backcolour},
    commentstyle=\color{codegreen},
    keywordstyle=\color{magenta},
    numberstyle=\tiny\color{codegray},
    stringstyle=\color{codepurple},
    basicstyle=\ttfamily\footnotesize,
    breakatwhitespace=false,
    breaklines=true,
    captionpos=b,
    keepspaces=true,
    numbers=left,
    numbersep=5pt,
    showspaces=false,
    showstringspaces=false,
    showtabs=false,
    tabsize=2,
    frame=single
}
\lstset{style=mystyle}

\hypersetup{
    colorlinks=true,
    linkcolor=blue,
    filecolor=magenta,
    urlcolor=cyan,
    pdftitle={KVRM-OS Implementation Guide},
    pdfauthor={Bobby Price}
}

\pagestyle{fancy}
\fancyhf{}
\rhead{KVRM-OS Implementation Guide v0.2}
\lhead{blackWeb Research}
\rfoot{Page \thepage}

\title{\textbf{KVRM-OS Implementation Guide}\\
\large A Step-by-Step Manual for Building KVRM-Vector\\
\vspace{0.5cm}
\normalsize Version 0.2}
\author{Bobby Price\\
blackWeb Research\\
\href{mailto:contact@blackweb.dev}{contact@blackweb.dev}}
\date{December 2025}

\begin{document}

\maketitle

\begin{abstract}
This document provides comprehensive step-by-step implementation instructions for building the KVRM-OS prototype using Python and modern ML tools. It covers project setup, data generation, micro-LLM training with Unsloth/LoRA, registry implementation, orchestration, testing, benchmarking, and deployment. By following this guide, you will produce a working KVRM-Vector prototype that demonstrates the key-value response mapping paradigm for model-native computing.

\textbf{Target Audience:} ML engineers, systems programmers, and researchers interested in neuro-symbolic AI and model-native computing.

\textbf{Prerequisites:} Python 3.10+, basic ML/PyTorch knowledge, GPU recommended (CPU fallback provided).
\end{abstract}

\tableofcontents
\newpage

\section{Overview and Prerequisites}

\subsection{What We're Building}

KVRM-Vector is a vector (dynamic array) data structure implemented entirely through micro-language models. Instead of traditional code like \texttt{list.append(x)}, operations are handled by fine-tuned LLMs that emit verified keys:

\begin{lstlisting}[language=Python, caption={Traditional vs KVRM Approach}]
# Traditional Python
my_list = []
my_list.append(42)
print(my_list[0])

# KVRM-Vector (semantic programming)
vec = KVRMVector()
vec.push(42)       # orchestrator -> push_llm -> {"action": "vec:push", "val": 42}
print(vec.get(0))  # orchestrator -> get_llm -> {"action": "vec:get", "idx": 0}
\end{lstlisting}

\subsection{System Architecture}

\begin{verbatim}
+------------------+     +------------------+     +------------------+
|   User Goal      | --> |   Orchestrator   | --> |   Micro-LLM      |
| "add 42 to vec"  |     |   (routes goal)  |     |   (emits key)    |
+------------------+     +------------------+     +------------------+
                                                          |
                                                          v
+------------------+     +------------------+     +------------------+
|   Result State   | <-- |   Executor       | <-- |   Registry       |
| vec = [42]       |     |   (runs action)  |     |   (key->action)  |
+------------------+     +------------------+     +------------------+
\end{verbatim}

\subsection{Components to Build}

\begin{enumerate}
    \item \textbf{Data Generator}: Creates (context $\to$ key) training pairs
    \item \textbf{Micro-LLMs}: 7 fine-tuned models (orchestrator + 6 operations)
    \item \textbf{Registry}: Immutable key $\to$ action mapping (SQLite-backed)
    \item \textbf{Executor}: Verified action execution engine
    \item \textbf{Orchestrator}: Goal routing and chaining logic
    \item \textbf{KVRMVector}: Main API class with list-like interface
    \item \textbf{Demo}: Interactive Gradio application
    \item \textbf{Benchmarks}: Performance comparison vs native Python
\end{enumerate}

\subsection{Prerequisites Checklist}

\begin{table}[h]
\centering
\begin{tabular}{lll}
\toprule
\textbf{Requirement} & \textbf{Minimum} & \textbf{Recommended} \\
\midrule
Python & 3.10 & 3.11+ \\
RAM & 16 GB & 32 GB \\
GPU & CPU fallback & CUDA/Apple Silicon \\
Disk Space & 10 GB & 50 GB \\
Training Time & 8-12h (CPU) & 2-4h (GPU) \\
\bottomrule
\end{tabular}
\caption{System Requirements}
\end{table}

\subsection{Estimated Timeline}

\begin{itemize}
    \item \textbf{Setup}: 15 minutes
    \item \textbf{Data Generation}: 30 minutes
    \item \textbf{Training}: 2-4 hours (GPU) / 8-12 hours (CPU)
    \item \textbf{Integration}: 1-2 hours
    \item \textbf{Testing \& Benchmarks}: 1 hour
    \item \textbf{Total}: 4-8 hours
\end{itemize}

\section{Project Setup}

\subsection{Directory Structure}

Create the following structure:

\begin{lstlisting}[language=bash, caption={Project Structure}]
kvrm-vector/
|-- README.md
|-- requirements.txt
|-- pyproject.toml
|-- src/
|   |-- __init__.py
|   |-- registry.py
|   |-- data_generator.py
|   |-- micro_llm.py
|   |-- orchestrator.py
|   |-- executor.py
|   |-- kvrm_vector.py
|-- training/
|   |-- train_micro.py
|   |-- configs/
|       |-- orchestrator.yaml
|       |-- push.yaml
|       |-- pop.yaml
|       |-- get.yaml
|       |-- sort.yaml
|       |-- create.yaml
|       |-- verify.yaml
|-- data/
|-- models/
|-- tests/
|   |-- test_registry.py
|   |-- test_micro_llm.py
|   |-- test_kvrm_vector.py
|   |-- test_integration.py
|-- benchmarks/
|   |-- benchmark_vs_native.py
|-- demo/
|   |-- gradio_app.py
|-- docs/
    |-- latex/
\end{lstlisting}

\subsection{Dependencies}

\begin{lstlisting}[caption={requirements.txt}]
# Core ML
torch>=2.0.0
transformers>=4.35.0
datasets>=2.14.0
peft>=0.6.0
accelerate>=0.24.0
bitsandbytes>=0.41.0

# Fast Training
unsloth>=2024.8

# Data Generation
faker>=18.0.0
pyyaml>=6.0

# Demo & Testing
gradio>=4.0.0
pytest>=7.4.0
pytest-cov>=4.1.0

# Utilities
tqdm>=4.65.0
numpy>=1.24.0
pandas>=2.0.0
\end{lstlisting}

\subsection{Installation Commands}

\begin{lstlisting}[language=bash, caption={Setup Commands}]
# Create project directory
mkdir -p kvrm-vector/{src,training/configs,data,models,tests,benchmarks,demo,docs}

# Initialize virtual environment
cd kvrm-vector
python -m venv venv
source venv/bin/activate  # On Windows: venv\Scripts\activate

# Install dependencies
pip install -r requirements.txt

# Initialize git
git init
echo "venv/\n*.pyc\n__pycache__/\nmodels/\ndata/*.jsonl" > .gitignore
git add .
git commit -m "Initial project setup"
\end{lstlisting}

\subsection{Verify Setup}

\begin{lstlisting}[language=Python, caption={verify\_setup.py}]
#!/usr/bin/env python3
"""Verify KVRM-Vector setup is complete."""
import sys
import psutil

def check_setup():
    print("=== KVRM-Vector Setup Verification ===\n")

    # Python version
    py_ver = sys.version_info
    print(f"Python: {py_ver.major}.{py_ver.minor}.{py_ver.micro}")
    assert py_ver >= (3, 10), "Python 3.10+ required"

    # PyTorch and GPU
    import torch
    print(f"PyTorch: {torch.__version__}")
    print(f"CUDA available: {torch.cuda.is_available()}")
    print(f"MPS available: {torch.backends.mps.is_available()}")

    # Memory
    ram_gb = psutil.virtual_memory().total / (1024**3)
    print(f"RAM: {ram_gb:.1f} GB")
    assert ram_gb >= 16, "16GB+ RAM recommended"

    # Key packages
    import transformers
    import peft
    import gradio
    import faker
    print(f"Transformers: {transformers.__version__}")
    print(f"PEFT: {peft.__version__}")
    print(f"Gradio: {gradio.__version__}")

    print("\n=== All checks passed! ===")

if __name__ == "__main__":
    check_setup()
\end{lstlisting}

\section{Registry Implementation}

The registry is the safety foundation of KVRM: an immutable mapping from keys to verified actions.

\subsection{Registry Design Principles}

\begin{enumerate}
    \item \textbf{Immutability}: Once frozen, no new keys can be added
    \item \textbf{Integrity}: Merkle hash verification prevents tampering
    \item \textbf{Determinism}: Same key always maps to same action
    \item \textbf{Completeness}: All valid keys are pre-registered
\end{enumerate}

\subsection{Core Registry Implementation}

\begin{lstlisting}[language=Python, caption={registry.py Core}]
import sqlite3
import hashlib
import json
from typing import Callable, Dict, Any, Optional
from dataclasses import dataclass
from functools import lru_cache

class InvalidKeyError(Exception):
    """Raised when a key is not in the registry."""
    pass

class RegistryFrozenError(Exception):
    """Raised when trying to modify a frozen registry."""
    pass

@dataclass
class RegistryEntry:
    key: str
    action_id: str
    description: str
    category: str

class KVRMRegistry:
    """Immutable key-to-action mapping with integrity verification."""

    def __init__(self, db_path: str = ":memory:"):
        self.db_path = db_path
        self.conn = sqlite3.connect(db_path)
        self.frozen = False
        self.integrity_hash: Optional[str] = None
        self.primitives: Dict[str, Callable] = {}
        self._create_schema()
        self._register_vector_primitives()

    def _create_schema(self):
        self.conn.execute("""
            CREATE TABLE IF NOT EXISTS registry (
                key TEXT PRIMARY KEY,
                action_id TEXT NOT NULL,
                description TEXT,
                category TEXT,
                created_at TIMESTAMP DEFAULT CURRENT_TIMESTAMP
            )
        """)
        self.conn.commit()

    def register(self, key: str, action: Callable,
                 description: str = "", category: str = "general"):
        if self.frozen:
            raise RegistryFrozenError("Cannot modify frozen registry")

        action_id = f"action_{id(action)}"
        self.primitives[action_id] = action

        self.conn.execute(
            "INSERT OR REPLACE INTO registry VALUES (?, ?, ?, ?, CURRENT_TIMESTAMP)",
            (key, action_id, description, category)
        )
        self.conn.commit()

    def freeze(self) -> str:
        """Freeze registry and compute integrity hash."""
        self.frozen = True
        self.integrity_hash = self._compute_hash()
        return self.integrity_hash

    def _compute_hash(self) -> str:
        rows = self.conn.execute(
            "SELECT key, action_id FROM registry ORDER BY key"
        ).fetchall()
        data = json.dumps(rows, sort_keys=True)
        return hashlib.sha256(data.encode()).hexdigest()

    @lru_cache(maxsize=1024)
    def lookup(self, key: str) -> Callable:
        row = self.conn.execute(
            "SELECT action_id FROM registry WHERE key = ?", (key,)
        ).fetchone()
        if row is None:
            raise InvalidKeyError(f"Key not found: {key}")
        return self.primitives[row[0]]

    def execute(self, key: str, state: Dict[str, Any],
                **kwargs) -> Dict[str, Any]:
        action = self.lookup(key)
        return action(state, **kwargs)
\end{lstlisting}

\subsection{Vector Primitives Registration}

\begin{lstlisting}[language=Python, caption={Vector Primitives}]
def _register_vector_primitives(self):
    """Register all vector operation primitives."""

    # Create operations
    self.register(
        "vec:create",
        lambda state, name: {**state, name: []},
        "Create new empty vector", "create"
    )
    self.register(
        "vec:create_with_capacity",
        lambda state, name, capacity: {**state, name: [], f"{name}_cap": capacity},
        "Create vector with capacity hint", "create"
    )

    # Push operations
    self.register(
        "vec:push",
        lambda state, name, value: {**state, name: state[name] + [value]},
        "Push value to end", "push"
    )
    self.register(
        "vec:push_front",
        lambda state, name, value: {**state, name: [value] + state[name]},
        "Push value to front", "push"
    )

    # Pop operations
    self.register(
        "vec:pop",
        lambda state, name: ({**state, name: state[name][:-1]}, state[name][-1]),
        "Pop from end", "pop"
    )
    self.register(
        "vec:pop_front",
        lambda state, name: ({**state, name: state[name][1:]}, state[name][0]),
        "Pop from front", "pop"
    )

    # Get operations
    self.register(
        "vec:get",
        lambda state, name, idx: state[name][idx],
        "Get value at index", "get"
    )
    self.register(
        "vec:length",
        lambda state, name: len(state[name]),
        "Get vector length", "get"
    )

    # Sort operations
    self.register(
        "vec:sort_asc",
        lambda state, name: {**state, name: sorted(state[name])},
        "Sort ascending", "sort"
    )
    self.register(
        "vec:sort_desc",
        lambda state, name: {**state, name: sorted(state[name], reverse=True)},
        "Sort descending", "sort"
    )
\end{lstlisting}

\section{Data Generation}

Training data consists of (context, key) pairs for each micro-LLM.

\subsection{Dataset Sizes}

\begin{table}[h]
\centering
\begin{tabular}{lrl}
\toprule
\textbf{Model} & \textbf{Examples} & \textbf{Description} \\
\midrule
orchestrator & 10,000 & Routes goals to operations \\
create & 3,000 & Vector creation variants \\
push & 4,000 & Push operations \\
pop & 3,000 & Pop operations \\
get & 3,000 & Access operations \\
sort & 3,000 & Sorting operations \\
verify & 3,000 & Validation operations \\
\midrule
\textbf{Total} & \textbf{29,000} & \\
\bottomrule
\end{tabular}
\caption{Training Dataset Sizes}
\end{table}

\subsection{Data Generation with Faker}

\begin{lstlisting}[language=Python, caption={Data Generator Core}]
from faker import Faker
import json
import random
from typing import List, Dict

fake = Faker()

def generate_push_examples(n: int = 4000) -> List[Dict]:
    """Generate push operation training examples."""
    examples = []

    templates = [
        "add {value} to {vec_name}",
        "push {value} onto {vec_name}",
        "append {value} to the {vec_name} vector",
        "insert {value} at the end of {vec_name}",
        "{vec_name}.push({value})",
        "put {value} in {vec_name}",
    ]

    for _ in range(n):
        vec_name = fake.word() + "_vec"
        value_type = random.choice(["int", "float", "string"])

        if value_type == "int":
            value = random.randint(-1000, 1000)
        elif value_type == "float":
            value = round(random.uniform(-100, 100), 2)
        else:
            value = f'"{fake.word()}"'

        template = random.choice(templates)
        context = template.format(value=value, vec_name=vec_name)

        key = {
            "action": "vec:push",
            "target": vec_name,
            "value": value
        }

        examples.append({
            "context": context,
            "key": json.dumps(key),
            "category": "push"
        })

    return examples
\end{lstlisting}

\subsection{Running Data Generation}

\begin{lstlisting}[language=bash, caption={Generate All Data}]
python src/data_generator.py --all --output-dir data/

# Expected output:
# Generated 10000 orchestrator examples -> data/orchestrator.jsonl
# Generated 3000 create examples -> data/create.jsonl
# Generated 4000 push examples -> data/push.jsonl
# Generated 3000 pop examples -> data/pop.jsonl
# Generated 3000 get examples -> data/get.jsonl
# Generated 3000 sort examples -> data/sort.jsonl
# Generated 3000 verify examples -> data/verify.jsonl
# Total: 29000 examples
\end{lstlisting}

\section{Micro-LLM Training}

\subsection{Training Configuration}

\begin{lstlisting}[language=yaml, caption={configs/push.yaml}]
model:
  base: "unsloth/Llama-3.2-1B"
  max_seq_length: 512
  load_in_4bit: true

lora:
  r: 16
  lora_alpha: 32
  lora_dropout: 0.05
  target_modules: ["q_proj", "k_proj", "v_proj", "o_proj"]

training:
  epochs: 3
  batch_size: 8
  gradient_accumulation: 4
  learning_rate: 2e-4
  warmup_ratio: 0.1
  weight_decay: 0.01
  fp16: true

data:
  path: "data/push.jsonl"
  train_split: 0.9

output:
  path: "models/push_llm"
  save_steps: 500
\end{lstlisting}

\subsection{Training Script}

\begin{lstlisting}[language=Python, caption={train\_micro.py}]
#!/usr/bin/env python3
"""Train all KVRM micro-LLMs using Unsloth/LoRA."""

from unsloth import FastLanguageModel
from datasets import load_dataset
from trl import SFTTrainer
from transformers import TrainingArguments
import yaml
import argparse

def train_model(config_path: str):
    with open(config_path) as f:
        cfg = yaml.safe_load(f)

    # Load base model with 4-bit quantization
    model, tokenizer = FastLanguageModel.from_pretrained(
        model_name=cfg["model"]["base"],
        max_seq_length=cfg["model"]["max_seq_length"],
        load_in_4bit=cfg["model"]["load_in_4bit"],
    )

    # Add LoRA adapters
    model = FastLanguageModel.get_peft_model(
        model,
        r=cfg["lora"]["r"],
        lora_alpha=cfg["lora"]["lora_alpha"],
        lora_dropout=cfg["lora"]["lora_dropout"],
        target_modules=cfg["lora"]["target_modules"],
    )

    # Load and format dataset
    dataset = load_dataset("json", data_files=cfg["data"]["path"])

    def format_prompt(example):
        return f"### Context:\n{example['context']}\n\n### Key:\n{example['key']}"

    # Training
    trainer = SFTTrainer(
        model=model,
        tokenizer=tokenizer,
        train_dataset=dataset["train"],
        dataset_text_field="text",
        max_seq_length=cfg["model"]["max_seq_length"],
        args=TrainingArguments(
            output_dir=cfg["output"]["path"],
            num_train_epochs=cfg["training"]["epochs"],
            per_device_train_batch_size=cfg["training"]["batch_size"],
            learning_rate=cfg["training"]["learning_rate"],
            fp16=cfg["training"]["fp16"],
            save_steps=cfg["output"]["save_steps"],
        ),
    )

    trainer.train()
    model.save_pretrained(cfg["output"]["path"])

if __name__ == "__main__":
    parser = argparse.ArgumentParser()
    parser.add_argument("--config", required=True)
    parser.add_argument("--all", action="store_true")
    args = parser.parse_args()

    if args.all:
        models = ["orchestrator", "create", "push", "pop", "get", "sort", "verify"]
        for m in models:
            train_model(f"training/configs/{m}.yaml")
    else:
        train_model(args.config)
\end{lstlisting}

\section{Micro-LLM Wrapper}

\begin{lstlisting}[language=Python, caption={micro\_llm.py}]
from functools import lru_cache
from transformers import AutoModelForCausalLM, AutoTokenizer
import torch
import json
import re

class MicroLLM:
    """Wrapper for fine-tuned micro-LLMs with caching and validation."""

    def __init__(self, model_path: str, valid_keys: set = None):
        self.model_path = model_path
        self.valid_keys = valid_keys or set()
        self.model = None
        self.tokenizer = None
        self._loaded = False

    def load(self):
        if self._loaded:
            return
        self.tokenizer = AutoTokenizer.from_pretrained(self.model_path)
        self.model = AutoModelForCausalLM.from_pretrained(
            self.model_path,
            torch_dtype=torch.float16,
            device_map="auto"
        )
        self._loaded = True

    @lru_cache(maxsize=10000)
    def infer(self, context: str) -> dict:
        """Infer key from context with caching."""
        self.load()

        prompt = f"### Context:\n{context}\n\n### Key:\n"
        inputs = self.tokenizer(prompt, return_tensors="pt").to(self.model.device)

        outputs = self.model.generate(
            **inputs,
            max_new_tokens=128,
            temperature=0.1,
            do_sample=True,
            pad_token_id=self.tokenizer.eos_token_id
        )

        response = self.tokenizer.decode(outputs[0], skip_special_tokens=True)
        key_text = response.split("### Key:\n")[-1].strip()

        # Parse and validate JSON key
        try:
            key = json.loads(key_text)
            if self.valid_keys and key.get("action") not in self.valid_keys:
                raise ValueError(f"Invalid action: {key.get('action')}")
            return key
        except json.JSONDecodeError:
            return {"error": "invalid_json", "raw": key_text}
\end{lstlisting}

\section{Orchestrator}

\begin{lstlisting}[language=Python, caption={orchestrator.py}]
from typing import Dict, Any, Optional
from .micro_llm import MicroLLM
from .registry import KVRMRegistry
from .executor import KVRMExecutor

class KVRMOrchestrator:
    """Routes high-level goals to micro-LLMs and chains execution."""

    def __init__(self, models_dir: str = "models"):
        self.models_dir = models_dir
        self.registry = KVRMRegistry()
        self.executor = KVRMExecutor(self.registry)

        # Load micro-LLMs
        self.orchestrator_llm = MicroLLM(f"{models_dir}/orchestrator_llm")
        self.op_llms = {
            "create": MicroLLM(f"{models_dir}/create_llm"),
            "push": MicroLLM(f"{models_dir}/push_llm"),
            "pop": MicroLLM(f"{models_dir}/pop_llm"),
            "get": MicroLLM(f"{models_dir}/get_llm"),
            "sort": MicroLLM(f"{models_dir}/sort_llm"),
            "verify": MicroLLM(f"{models_dir}/verify_llm"),
        }

        self.state: Dict[str, Any] = {}
        self.trace: list = []

    def process_goal(self, goal: str) -> Dict[str, Any]:
        """Process a natural language goal through the KVRM pipeline."""
        self.trace = []

        # Step 1: Route goal to operation
        route_key = self.orchestrator_llm.infer(goal)
        self.trace.append({"step": "route", "key": route_key})

        if "error" in route_key:
            return {"success": False, "error": route_key["error"]}

        # Step 2: Get operation-specific key
        op_type = route_key.get("operation")
        if op_type not in self.op_llms:
            return {"success": False, "error": f"Unknown operation: {op_type}"}

        op_key = self.op_llms[op_type].infer(goal)
        self.trace.append({"step": "operation", "key": op_key})

        if "error" in op_key:
            return {"success": False, "error": op_key["error"]}

        # Step 3: Execute through registry
        try:
            result = self.executor.execute(op_key, self.state)
            self.trace.append({"step": "execute", "result": result})
            return {"success": True, "result": result, "trace": self.trace}
        except Exception as e:
            return {"success": False, "error": str(e), "trace": self.trace}
\end{lstlisting}

\section{Executor}

\begin{lstlisting}[language=Python, caption={executor.py}]
from typing import Dict, Any
from .registry import KVRMRegistry, InvalidKeyError

class KVRMExecutor:
    """Executes verified actions from the registry."""

    def __init__(self, registry: KVRMRegistry):
        self.registry = registry

    def execute(self, key_data: Dict[str, Any],
                state: Dict[str, Any]) -> Dict[str, Any]:
        """Execute a key against the current state."""
        action_key = key_data.get("action")
        if not action_key:
            raise ValueError("No action specified in key")

        # Extract parameters
        params = {k: v for k, v in key_data.items() if k != "action"}

        # Lookup and execute
        action = self.registry.lookup(action_key)
        new_state = action(state, **params)

        return new_state
\end{lstlisting}

\section{KVRMVector Main Class}

\begin{lstlisting}[language=Python, caption={kvrm\_vector.py}]
from typing import Any, Iterator
from .orchestrator import KVRMOrchestrator

class KVRMVector:
    """A vector data structure implemented via micro-LLMs."""

    def __init__(self, name: str = "default_vec"):
        self.name = name
        self.orchestrator = KVRMOrchestrator()
        self._initialize()

    def _initialize(self):
        self.orchestrator.process_goal(f"create a new vector called {self.name}")

    def push(self, value: Any) -> None:
        self.orchestrator.process_goal(f"add {value} to {self.name}")

    def pop(self) -> Any:
        result = self.orchestrator.process_goal(f"remove last element from {self.name}")
        return result.get("result")

    def get(self, index: int) -> Any:
        result = self.orchestrator.process_goal(f"get element {index} from {self.name}")
        return result.get("result")

    def __len__(self) -> int:
        result = self.orchestrator.process_goal(f"get length of {self.name}")
        return result.get("result", 0)

    def sort(self, reverse: bool = False) -> None:
        direction = "descending" if reverse else "ascending"
        self.orchestrator.process_goal(f"sort {self.name} in {direction} order")

    def __iter__(self) -> Iterator:
        for i in range(len(self)):
            yield self.get(i)

    def __repr__(self) -> str:
        return f"KVRMVector({self.name}, {list(self)})"
\end{lstlisting}

\section{Testing}

\begin{lstlisting}[language=Python, caption={test\_kvrm\_vector.py}]
import pytest
from src.kvrm_vector import KVRMVector

class TestKVRMVector:
    def test_create_empty(self):
        vec = KVRMVector("test")
        assert len(vec) == 0

    def test_push_single(self):
        vec = KVRMVector("test")
        vec.push(42)
        assert len(vec) == 1
        assert vec.get(0) == 42

    def test_push_multiple(self):
        vec = KVRMVector("test")
        for i in range(5):
            vec.push(i)
        assert len(vec) == 5
        assert list(vec) == [0, 1, 2, 3, 4]

    def test_pop(self):
        vec = KVRMVector("test")
        vec.push(1)
        vec.push(2)
        val = vec.pop()
        assert val == 2
        assert len(vec) == 1

    def test_sort(self):
        vec = KVRMVector("test")
        for v in [3, 1, 4, 1, 5]:
            vec.push(v)
        vec.sort()
        assert list(vec) == [1, 1, 3, 4, 5]
\end{lstlisting}

\section{Benchmarking}

\begin{lstlisting}[language=Python, caption={benchmark\_vs\_native.py}]
import time
from src.kvrm_vector import KVRMVector

def benchmark_push(n: int = 1000):
    # Native Python
    native_list = []
    start = time.perf_counter()
    for i in range(n):
        native_list.append(i)
    native_time = time.perf_counter() - start

    # KVRM Vector
    kvrm_vec = KVRMVector("bench")
    start = time.perf_counter()
    for i in range(n):
        kvrm_vec.push(i)
    kvrm_time = time.perf_counter() - start

    print(f"Push {n} elements:")
    print(f"  Native: {native_time*1000:.2f}ms")
    print(f"  KVRM:   {kvrm_time*1000:.2f}ms")
    print(f"  Overhead: {kvrm_time/native_time:.1f}x")

if __name__ == "__main__":
    benchmark_push(1000)
\end{lstlisting}

\section{Demo Application}

\begin{lstlisting}[language=Python, caption={gradio\_app.py}]
import gradio as gr
from src.orchestrator import KVRMOrchestrator

orch = KVRMOrchestrator()

def process_command(command: str, history: list):
    result = orch.process_goal(command)
    history.append({"command": command, "result": result})
    state_display = str(orch.state)
    trace_display = str(result.get("trace", []))
    return state_display, trace_display, history

with gr.Blocks(title="KVRM-Vector Demo") as demo:
    gr.Markdown("# KVRM-Vector Demo")
    gr.Markdown("Enter natural language commands to manipulate vectors.")

    with gr.Row():
        command_input = gr.Textbox(label="Command", placeholder="create a vector called nums")
        submit_btn = gr.Button("Execute")

    with gr.Row():
        state_output = gr.Textbox(label="Current State", lines=5)
        trace_output = gr.Textbox(label="Execution Trace", lines=5)

    history_state = gr.State([])

    submit_btn.click(
        process_command,
        inputs=[command_input, history_state],
        outputs=[state_output, trace_output, history_state]
    )

if __name__ == "__main__":
    demo.launch()
\end{lstlisting}

\section{Troubleshooting}

\begin{table}[h]
\centering
\begin{tabular}{p{4cm}p{4cm}p{4cm}}
\toprule
\textbf{Issue} & \textbf{Cause} & \textbf{Solution} \\
\midrule
Out of memory & Model too large & Use 4-bit quantization \\
Invalid JSON output & Undertrained model & Increase epochs, add more data \\
Wrong operation routed & Data imbalance & Balance examples per class \\
High latency & No caching & Enable LRU cache \\
Overfitting & Too many epochs & Add early stopping \\
\bottomrule
\end{tabular}
\caption{Common Issues and Solutions}
\end{table}

\section{Deployment}

\subsection{Hugging Face Spaces}

\begin{lstlisting}[language=bash, caption={Deploy to HF Spaces}]
# Login to Hugging Face
huggingface-cli login

# Push models
huggingface-cli upload kvrm-vector-demo ./models --repo-type space

# Create Space with Gradio
# Add app.py and requirements.txt to the Space
\end{lstlisting}

\part{KVRM-CPU: Semantic Instruction Decode Implementation}

\section{KVRM-CPU Overview}

KVRM-CPU extends the KVRM paradigm to the lowest level of computing: CPU instruction decode. Instead of traditional hardcoded opcode parsing, a semantic micro-LLM interprets instructions and emits verified registry keys.

\subsection{The CPU Paradigm Shift}

\begin{lstlisting}[language=Python, caption={Traditional vs KVRM CPU Decode}]
# Traditional CPU (hardcoded silicon)
fetch -> decode -> execute
# Each opcode has fixed binary mapping

# KVRM-CPU (semantic understanding)
fetch -> decode_llm -> key -> verified_execute
# Natural language or assembly -> semantic key -> verified primitive
\end{lstlisting}

\subsection{Why CPU-Level KVRM Matters}

\begin{enumerate}
    \item \textbf{Proves Universality}: If KVRM works at the CPU level, it works everywhere
    \item \textbf{Minimal Surface}: Only 16 primitives cover a complete instruction set
    \item \textbf{Full Auditability}: Every instruction decoded and executed is logged
    \item \textbf{Semantic Flexibility}: "increment R0" and "R0++" decode identically
\end{enumerate}

\section{KVRM-CPU Project Setup}

\subsection{Directory Structure}

\begin{lstlisting}[language=bash, caption={KVRM-CPU Project Structure}]
kvrm-cpu/
|-- README.md
|-- requirements.txt
|-- src/
|   |-- kvrm_cpu/
|       |-- __init__.py
|       |-- state.py        # CPUState dataclass
|       |-- registry.py     # 16 CPU primitives
|       |-- decode_llm.py   # Semantic decoder
|       |-- cpu.py          # Main orchestrator
|-- programs/
|   |-- sum_1_to_10.asm
|   |-- fibonacci.asm
|   |-- countdown.asm
|   |-- countup_to_negative.asm
|-- training/
|   |-- generate_cpu_data.py
|   |-- train_decode.py
|-- tests/
|   |-- test_state.py
|   |-- test_registry.py
|   |-- test_decode.py
|   |-- test_programs.py
|-- benchmarks/
|-- models/
\end{lstlisting}

\section{CPUState Implementation}

The CPUState is an immutable dataclass representing the complete CPU state.

\begin{lstlisting}[language=Python, caption={state.py - CPUState Definition}]
from dataclasses import dataclass, field
from typing import Dict, List, Optional

@dataclass(frozen=False)
class CPUState:
    """Immutable CPU state with registers, flags, and memory."""

    registers: Dict[str, int] = field(default_factory=lambda: {
        f"R{i}": 0 for i in range(8)  # R0-R7
    })
    pc: int = 0                       # Program counter
    flags: Dict[str, bool] = field(default_factory=lambda: {
        "ZF": False,  # Zero flag
        "SF": False,  # Sign flag (negative)
    })
    memory: List[str] = field(default_factory=list)  # Instructions
    halted: bool = False
    cycle_count: int = 0

    def set_register(self, reg: str, value: int) -> 'CPUState':
        """Return new state with register updated."""
        new_regs = dict(self.registers)
        new_regs[reg] = value
        return CPUState(
            registers=new_regs,
            pc=self.pc,
            flags=dict(self.flags),
            memory=self.memory,
            halted=self.halted,
            cycle_count=self.cycle_count
        )

    def snapshot(self) -> dict:
        """Return serializable state snapshot for tracing."""
        return {
            "registers": dict(self.registers),
            "pc": self.pc,
            "flags": dict(self.flags),
            "halted": self.halted,
            "cycle_count": self.cycle_count
        }
\end{lstlisting}

\section{CPU Registry: 16 Verified Primitives}

The registry contains all valid CPU operations. Once frozen, no new operations can be added.

\subsection{Instruction Set Architecture}

\begin{table}[h]
\centering
\begin{tabular}{llp{6cm}}
\toprule
\textbf{Key} & \textbf{Params} & \textbf{Description} \\
\midrule
\texttt{OP\_MOV\_REG\_IMM} & dest, value & Load immediate into register \\
\texttt{OP\_MOV\_REG\_REG} & dest, src & Register-to-register copy \\
\texttt{OP\_ADD} & dest, src1, src2 & Addition (dest = src1 + src2) \\
\texttt{OP\_SUB} & dest, src1, src2 & Subtraction (dest = src1 - src2) \\
\texttt{OP\_MUL} & dest, src1, src2 & Multiplication \\
\texttt{OP\_INC} & dest & Increment by 1, sets flags \\
\texttt{OP\_DEC} & dest & Decrement by 1, sets flags \\
\texttt{OP\_CMP} & src1, src2 & Compare (sets ZF, SF) \\
\texttt{OP\_JMP} & addr & Unconditional jump \\
\texttt{OP\_JZ} & addr & Jump if zero (ZF set) \\
\texttt{OP\_JNZ} & addr & Jump if not zero \\
\texttt{OP\_JS} & addr & Jump if sign (SF set, negative) \\
\texttt{OP\_JNS} & addr & Jump if not sign \\
\texttt{OP\_HALT} & - & Stop execution \\
\texttt{OP\_NOP} & - & No operation \\
\texttt{OP\_INVALID} & raw & Error handler \\
\bottomrule
\end{tabular}
\caption{KVRM-CPU Instruction Set (16 Primitives)}
\end{table}

\subsection{Registry Implementation}

\begin{lstlisting}[language=Python, caption={registry.py - CPU Primitives}]
class CPURegistry:
    """Frozen registry of CPU operation primitives."""

    def __init__(self):
        self._primitives: Dict[str, Callable] = {}
        self._frozen = False
        self._register_all_primitives()
        self.freeze()

    def _register_all_primitives(self):
        """Register all 16 CPU operations."""

        # Data Movement
        self.register("OP_MOV_REG_IMM", self._op_mov_reg_imm)
        self.register("OP_MOV_REG_REG", self._op_mov_reg_reg)

        # Arithmetic
        self.register("OP_ADD", self._op_add)
        self.register("OP_SUB", self._op_sub)
        self.register("OP_MUL", self._op_mul)
        self.register("OP_INC", self._op_inc)
        self.register("OP_DEC", self._op_dec)

        # Comparison
        self.register("OP_CMP", self._op_cmp)

        # Control Flow
        self.register("OP_JMP", self._op_jmp)
        self.register("OP_JZ", self._op_jz)
        self.register("OP_JNZ", self._op_jnz)
        self.register("OP_JS", self._op_js)
        self.register("OP_JNS", self._op_jns)

        # Special
        self.register("OP_HALT", self._op_halt)
        self.register("OP_NOP", self._op_nop)
        self.register("OP_INVALID", self._op_invalid)

    def _op_inc(self, state: CPUState, params: dict) -> CPUState:
        """Increment register by 1, set flags."""
        dest = params["dest"]
        value = state.registers[dest] + 1

        new_state = state.set_register(dest, value)
        new_state.flags["ZF"] = (value == 0)
        new_state.flags["SF"] = (value < 0)
        new_state.pc += 1
        new_state.cycle_count += 1
        return new_state

    def _op_js(self, state: CPUState, params: dict) -> CPUState:
        """Jump if sign flag set (negative result)."""
        if state.flags["SF"]:
            new_pc = params["addr"]
        else:
            new_pc = state.pc + 1

        new_state = CPUState(
            registers=dict(state.registers),
            pc=new_pc,
            flags=dict(state.flags),
            memory=state.memory,
            halted=state.halted,
            cycle_count=state.cycle_count + 1
        )
        return new_state
\end{lstlisting}

\section{DecodeLLM: Semantic Instruction Decoder}

The DecodeLLM converts assembly text to verified registry keys.

\subsection{Mock Mode (Rule-Based)}

For development without GPU, use regex-based parsing:

\begin{lstlisting}[language=Python, caption={decode\_llm.py - Mock Decoder}]
import re
from typing import Tuple, Dict, Optional

class DecodeLLM:
    """Semantic instruction decoder with mock and real modes."""

    def __init__(self, mock_mode: bool = True, model_path: str = None):
        self.mock_mode = mock_mode
        self.model_path = model_path
        self.model = None

        # Patterns for mock mode
        self.patterns = {
            r'^MOV\s+(\w+)\s*,\s*(-?\d+)$': self._decode_mov_imm,
            r'^MOV\s+(\w+)\s*,\s*(\w+)$': self._decode_mov_reg,
            r'^(ADD|SUB|MUL)\s+(\w+)\s*,\s*(\w+)\s*,\s*(\w+)$': self._decode_arith,
            r'^(INC|DEC)\s+(\w+)$': self._decode_inc_dec,
            r'^CMP\s+(\w+)\s*,\s*(\w+)$': self._decode_cmp,
            r'^(JMP|JZ|JNZ|JS|JNS)\s+(\w+)$': self._decode_jump,
            r'^HALT$': lambda m: ("OP_HALT", {}),
            r'^NOP$': lambda m: ("OP_NOP", {}),
        }

    def decode(self, instruction: str, labels: Dict[str, int] = None
              ) -> Tuple[str, dict]:
        """Decode instruction to (key, params) tuple."""
        labels = labels or {}
        instr = instruction.strip().upper()

        # Remove comments
        if ';' in instr:
            instr = instr.split(';')[0].strip()

        if not instr:
            return ("OP_NOP", {})

        if self.mock_mode:
            return self._mock_decode(instr, labels)
        else:
            return self._real_decode(instr, labels)

    def _mock_decode(self, instr: str, labels: Dict[str, int]
                    ) -> Tuple[str, dict]:
        """Rule-based decoding for mock mode."""
        for pattern, handler in self.patterns.items():
            match = re.match(pattern, instr, re.IGNORECASE)
            if match:
                result = handler(match)
                # Resolve labels in jump addresses
                if result[0].startswith("OP_J") and result[0] != "OP_JMP":
                    addr = result[1].get("addr")
                    if isinstance(addr, str) and addr in labels:
                        result[1]["addr"] = labels[addr]
                return result

        return ("OP_INVALID", {"raw": instr})

    def _decode_inc_dec(self, match) -> Tuple[str, dict]:
        """Decode INC/DEC instruction."""
        op = match.group(1).upper()
        dest = match.group(2).upper()
        key = f"OP_{op}"
        return (key, {"dest": dest})
\end{lstlisting}

\subsection{Real Mode (Trained Model)}

\begin{lstlisting}[language=Python, caption={decode\_llm.py - Real Mode}]
def _real_decode(self, instr: str, labels: Dict[str, int]
                ) -> Tuple[str, dict]:
    """Use trained micro-LLM for semantic decoding."""
    if self.model is None:
        self._load_model()

    # Construct prompt
    prompt = f"### Context:\n{instr}\n\n### Key:\n"

    inputs = self.tokenizer(prompt, return_tensors="pt")
    inputs = inputs.to(self.model.device)

    outputs = self.model.generate(
        **inputs,
        max_new_tokens=64,
        temperature=0.1,
        do_sample=True,
        pad_token_id=self.tokenizer.eos_token_id
    )

    response = self.tokenizer.decode(outputs[0], skip_special_tokens=True)
    key_text = response.split("### Key:\n")[-1].strip()

    try:
        result = json.loads(key_text)
        key = result.get("key", "OP_INVALID")
        params = result.get("params", {})

        # Resolve labels
        if "addr" in params and params["addr"] in labels:
            params["addr"] = labels[params["addr"]]

        return (key, params)
    except json.JSONDecodeError:
        return ("OP_INVALID", {"raw": instr, "parse_error": key_text})
\end{lstlisting}

\section{CPU Orchestrator}

The KVRMCPU class ties everything together: fetch, decode, execute.

\begin{lstlisting}[language=Python, caption={cpu.py - Main Orchestrator}]
class KVRMCPU:
    """KVRM-based CPU emulator with full execution tracing."""

    def __init__(self, mock_mode: bool = True):
        self.registry = get_registry()
        self.decoder = DecodeLLM(mock_mode=mock_mode)
        self.state: Optional[CPUState] = None
        self.labels: Dict[str, int] = {}
        self.trace: List[dict] = []

    def load_program(self, program: List[str]):
        """Parse labels and load instructions into memory."""
        instructions = []
        self.labels = {}

        for line in program:
            line = line.strip()
            if not line or line.startswith(';'):
                continue

            # Parse labels
            if ':' in line and not line.startswith(';'):
                parts = line.split(':', 1)
                label = parts[0].strip()
                self.labels[label] = len(instructions)
                line = parts[1].strip() if len(parts) > 1 else ''

            if line and not line.startswith(';'):
                instructions.append(line)

        self.state = create_initial_state(instructions)
        self.trace = []

    def step(self) -> dict:
        """Execute single instruction cycle."""
        if self.state.halted or self.state.pc >= len(self.state.memory):
            return {"halted": True}

        # Fetch
        instruction = self.state.memory[self.state.pc]

        # Decode (semantic!)
        key, params = self.decoder.decode(instruction, self.labels)

        # Capture pre-state
        pre_state = self.state.snapshot()

        # Execute via registry
        self.state = self.registry.execute(self.state, key, params)

        # Record trace
        entry = {
            "cycle": self.state.cycle_count,
            "instruction": instruction,
            "key": key,
            "params": params,
            "pre_state": pre_state,
            "post_state": self.state.snapshot()
        }
        self.trace.append(entry)

        return entry

    def run(self, max_cycles: int = 10000) -> List[dict]:
        """Run until HALT or max cycles."""
        while not self.state.halted and self.state.cycle_count < max_cycles:
            self.step()
        return self.trace
\end{lstlisting}

\section{Training Data Generation}

Generate 50,000 training samples for the decode micro-LLM.

\subsection{Dataset Composition}

\begin{table}[h]
\centering
\begin{tabular}{lrl}
\toprule
\textbf{Category} & \textbf{Count} & \textbf{Examples} \\
\midrule
MOV reg, imm & 10,000 & "MOV R3, 42", "mov r1, 0" \\
MOV reg, reg & 5,000 & "MOV R1, R2" \\
ADD/SUB/MUL & 12,000 & "ADD R3, R1, R2", "sub R0, R5, R4" \\
INC/DEC & 3,000 & "INC R0", "DEC R7" \\
CMP & 5,000 & "CMP R1, R2" \\
JMP/JZ/JNZ/JS/JNS & 10,000 & "JMP 5", "JZ loop", "JS negative" \\
HALT/NOP & 2,000 & "HALT", "NOP" \\
Invalid & 3,000 & "FOO R1", "ADD", "" \\
\midrule
\textbf{Total} & \textbf{50,000} & \\
\bottomrule
\end{tabular}
\caption{KVRM-CPU Training Dataset}
\end{table}

\subsection{Data Augmentation}

\begin{lstlisting}[language=Python, caption={generate\_cpu\_data.py}]
import random
import json

def generate_mov_imm_examples(n: int = 10000) -> list:
    """Generate MOV register, immediate examples."""
    examples = []
    registers = [f"R{i}" for i in range(8)]

    templates = [
        "MOV {reg}, {val}",
        "mov {reg}, {val}",
        "Mov {reg} {val}",      # Missing comma
        "{reg} = {val}",        # Assignment style
        "load {val} into {reg}",# Natural language
    ]

    for _ in range(n):
        reg = random.choice(registers)
        val = random.randint(-1000, 1000)
        template = random.choice(templates)

        instruction = template.format(reg=reg, val=val)

        key = {
            "key": "OP_MOV_REG_IMM",
            "params": {"dest": reg.upper(), "value": val}
        }

        examples.append({
            "instruction": instruction,
            "key": json.dumps(key)
        })

    return examples

def generate_inc_dec_examples(n: int = 3000) -> list:
    """Generate INC/DEC examples with natural language variants."""
    examples = []
    registers = [f"R{i}" for i in range(8)]

    inc_templates = [
        "INC {reg}",
        "inc {reg}",
        "{reg}++",
        "increment {reg}",
        "add 1 to {reg}",
    ]

    dec_templates = [
        "DEC {reg}",
        "dec {reg}",
        "{reg}--",
        "decrement {reg}",
        "subtract 1 from {reg}",
    ]

    for _ in range(n // 2):
        reg = random.choice(registers)

        # INC
        template = random.choice(inc_templates)
        examples.append({
            "instruction": template.format(reg=reg),
            "key": json.dumps({
                "key": "OP_INC",
                "params": {"dest": reg.upper()}
            })
        })

        # DEC
        template = random.choice(dec_templates)
        examples.append({
            "instruction": template.format(reg=reg),
            "key": json.dumps({
                "key": "OP_DEC",
                "params": {"dest": reg.upper()}
            })
        })

    return examples
\end{lstlisting}

\section{Example Programs}

\subsection{Sum 1 to 10}

\begin{lstlisting}[caption={programs/sum\_1\_to\_10.asm}]
; Sum 1 to 10 - Expected result: R0 = 55
    MOV R0, 0       ; sum = 0
    MOV R1, 1       ; counter = 1
    MOV R2, 11      ; limit
    MOV R3, 1       ; increment
loop:
    ADD R0, R0, R1  ; sum += counter
    ADD R1, R1, R3  ; counter++
    CMP R1, R2      ; compare to limit
    JNZ loop        ; continue if not equal
    HALT            ; R0 = 55
\end{lstlisting}

\subsection{Countdown with Sign Flag}

\begin{lstlisting}[caption={programs/countdown.asm}]
; Countdown from 5, using DEC and JS
; Demonstrates sign flag transitions
    MOV R0, 5       ; Start at 5
    MOV R1, 0       ; Decrement count
loop:
    DEC R0          ; counter-- (sets SF when negative)
    INC R1          ; count decrements
    JS done         ; Jump if negative (SF set)
    JMP loop        ; Continue if >= 0
done:
    HALT            ; R0 = -1, R1 = 6
\end{lstlisting}

\section{Testing KVRM-CPU}

\begin{lstlisting}[language=Python, caption={test\_programs.py}]
import pytest
from kvrm_cpu.cpu import KVRMCPU

class TestPrograms:
    def test_sum_1_to_10(self):
        """Sum 1 to 10 should produce R0 = 55."""
        cpu = KVRMCPU(mock_mode=True)

        program = [
            "MOV R0, 0",
            "MOV R1, 1",
            "MOV R2, 11",
            "MOV R3, 1",
            "loop: ADD R0, R0, R1",
            "ADD R1, R1, R3",
            "CMP R1, R2",
            "JNZ loop",
            "HALT"
        ]

        cpu.load_program(program)
        cpu.run()

        assert cpu.state.registers["R0"] == 55
        assert cpu.state.halted is True

    def test_countdown(self):
        """Countdown should demonstrate JS behavior."""
        cpu = KVRMCPU(mock_mode=True)

        program = [
            "MOV R0, 5",
            "MOV R1, 0",
            "loop: DEC R0",
            "INC R1",
            "JS done",
            "JMP loop",
            "done: HALT"
        ]

        cpu.load_program(program)
        cpu.run()

        assert cpu.state.registers["R0"] == -1  # Went negative
        assert cpu.state.registers["R1"] == 6   # 6 decrements
\end{lstlisting}

\section{KVRM-CPU Benchmarks}

\begin{lstlisting}[language=Python, caption={cpu\_benchmark.py}]
import time
from kvrm_cpu.cpu import KVRMCPU

def benchmark_sum_program(iterations: int = 100):
    """Benchmark sum 1 to 10 program."""
    program = [
        "MOV R0, 0", "MOV R1, 1", "MOV R2, 11", "MOV R3, 1",
        "loop: ADD R0, R0, R1", "ADD R1, R1, R3",
        "CMP R1, R2", "JNZ loop", "HALT"
    ]

    # KVRM Mock Mode
    cpu = KVRMCPU(mock_mode=True)
    start = time.perf_counter()
    for _ in range(iterations):
        cpu.load_program(program)
        cpu.run()
    mock_time = time.perf_counter() - start

    print(f"KVRM Mock: {mock_time*1000/iterations:.2f}ms per run")
    print(f"Result: R0 = {cpu.state.registers['R0']}")

    return {
        "mock_time_ms": mock_time * 1000 / iterations,
        "result": cpu.state.registers["R0"]
    }

if __name__ == "__main__":
    results = benchmark_sum_program()
    print(f"\nVerified: R0 = {results['result']} (expected 55)")
\end{lstlisting}

\section{Conclusion}

This guide provides comprehensive implementation instructions for both KVRM-Vector and KVRM-CPU. Together, these prototypes demonstrate:

\begin{itemize}
    \item \textbf{Key-only emission} eliminates hallucination by construction
    \item \textbf{Micro-LLMs} can replace traditional logic at any level (data structures to CPU decode)
    \item \textbf{Caching} reduces overhead to acceptable levels for practical use
    \item \textbf{Full auditability} via execution traces for every operation
    \item \textbf{Semantic flexibility} allowing natural language and formal notation interchangeably
    \item \textbf{Universality} proven from high-level vector operations down to CPU instruction decode
\end{itemize}

\subsection{Component Summary}

\begin{table}[h]
\centering
\begin{tabular}{llrl}
\toprule
\textbf{Component} & \textbf{Purpose} & \textbf{Primitives} & \textbf{Status} \\
\midrule
KVRM-Vector & Dynamic array operations & 9 & Complete \\
KVRM-CPU & CPU instruction decode & 16 & Complete \\
\midrule
\textbf{Total} & & \textbf{25} & \\
\bottomrule
\end{tabular}
\caption{KVRM-OS Implementation Components}
\end{table}

\textbf{Next Steps:}
\begin{enumerate}
    \item Extend to more data structures (HashMap, Tree, Graph)
    \item Add concurrency primitives (Lock, Semaphore, Channel)
    \item Implement file system primitives (Open, Read, Write, Close)
    \item Network stack primitives (Connect, Send, Receive)
    \item Implement self-improvement loop for model updates
    \item Target bare-metal deployment on custom hardware
\end{enumerate}

\textbf{Repository Links:}
\begin{itemize}
    \item KVRM-Vector: \url{https://github.com/robertcprice/kvrm-vector}
    \item KVRM-CPU: \url{https://github.com/robertcprice/kvrm-cpu}
\end{itemize}

\end{document}
